% Paper v15 appendix: data, training details and hyperparameters (moved out of the method), inference, CRF
% derivations, evaluation protocols, baselines, supplementary results.

\section{Data Construction and Statistics}
\label{app:data}

\begin{table}[t]
\centering\small
\resizebox{\columnwidth}{!}{\begin{tabular}{lrrrrrrr}
\toprule
Split & Mtg. & Windows & Hours & Mean\,s & Utt. & Multi\,\% & Ovl.\,\% \\
\midrule
\multicolumn{8}{l}{\emph{En$\rightarrow$Zh}} \\
Train & 137 & 21{,}581 & 36.8 & 6.1 & 54{,}824 & 38.5 & 38.2 \\
Dev & 18 & 3{,}203 & 5.2 & 5.9 & 8{,}206 & 40.0 & 39.6 \\
Test & 16 & 2{,}375 & 3.9 & 5.9 & 5{,}790 & 35.0 & 34.7 \\
\multicolumn{8}{l}{\emph{Zh$\rightarrow$En}} \\
Train & 127 & 27{,}953 & 48.5 & 6.3 & 69{,}540 & 37.4 & 34.6 \\
Dev & 15 & 3{,}366 & 5.7 & 6.1 & 8{,}340 & 39.3 & 36.6 \\
Test & 16 & 3{,}317 & 5.8 & 6.3 & 8{,}634 & 36.5 & 34.5 \\
\multicolumn{8}{l}{\emph{Zh$\rightarrow$En, AliMeeting (real meetings)}} \\
Test & 20 & 4{,}243 & 7.5 & 6.4 & 9{,}730 & 33.4 & 33.2 \\
\bottomrule
\end{tabular}}
\caption{Reference windows of the two directions. En$\rightarrow$Zh: AMI IHM-mix with Chinese references;
Zh$\rightarrow$En: the re-voiced Chinese meetings (Appendix~\ref{app:zhdata}) with the AMI English transcripts as
references. \emph{Multi}: windows with two or more speakers; \emph{Ovl.}: windows in which two utterances overlap.}
\label{tab:data}
\end{table}


\paragraph{Reference windows.} Utterances that overlap or touch are grouped into reference windows of at most 20\,s
(Table~\ref{tab:data}); training targets, text references and the overlap study use these windows, whereas test
decoding runs on whole meetings through online detector windows (Appendix~\ref{app:online}).

\paragraph{Translations.} Every AMI utterance (IHM segmentation, official split) is translated by three MT systems
(Hy-MT2-30B-A3B, NiuTrans LMT-60-4B, MiLMMT-46-4B); one candidate is selected by reference-free quality estimation
(XCOMET-XL~\citep{guerreiro2024xcomet} and an XLM-R-XL estimator) and chrF agreement, synthesised in the speaker's
voice and checked for audibility and intelligibility. A window enters the data only if all its non-hesitation
utterances have a translation whose speech decodes.

\paragraph{Chinese AMI meetings.}\label{app:zhdata} Every utterance's Chinese translation is re-spoken by the UniSS
Talker with one voice prompt per speaker, trimmed, set to the speaker's median loudness in the English recording,
placed at the utterance's original start and time-stretched towards the English duration within $[0.8,1.5]$; a
speaker's next clip waits until the previous one ends. Clips failing an intelligibility check (CER $>0.3$) are left out
of the meeting, and 13 training meetings whose diarization error stayed high were dropped (none in test).

\paragraph{AliMeeting.} AliMeeting is used for testing only. The main results use DiariST's human English
translations~\citep{yang2024diarist}; the overlap study (Tables~\ref{tab:overlap} and~\ref{tab:ovmain8}) uses machine
translations (Hy-MT2-30B) of the human transcripts, without a human check. The eight AliMeeting evaluation meetings
(\num{v11:ae:hours}\,h) of the locked test were never used for training, development or any earlier test.

\section{Training Implementation Details and Hyperparameters}
\label{app:hyper}
\label{app:train}

Table~\ref{tab:hyper} lists every value that the method leaves symbolic; the subsections below give the procedures.

\begin{table*}[t]
\centering\small
\setlength{\tabcolsep}{4pt}
\renewcommand{\arraystretch}{1.05}
\begin{tabular}{@{}l>{\raggedright\arraybackslash}p{4.6cm}l>{\raggedright\arraybackslash}p{4.6cm}@{}}
\toprule
\multicolumn{2}{@{}l}{\emph{Thinker and speaker tokens (\S\ref{sec:method-thinker})}} & \multicolumn{2}{l}{\emph{Talker and speaking rate (\S\ref{sec:method-queue})}} \\
Encoder / Thinker & Qwen3-Omni-30B-A3B-Instruct & Talker & UniSS, fine-tuned for block-wise speech \\
Hidden size $d$ & 2{,}048 & Voice prompt $\bm{v}_k$ & solo runs $\ge$0.4\,s, at most 8\,s \\
Acoustic token hop & 80\,ms & Prompt encoding & 32 BiCodec global tokens \\
Diarizer & Nemotron 3 Diarization & Speech codes & BiCodec semantic, one per 20\,ms \\
Slots $S$ / frame & 8 / 80\,ms (10\,ms posteriors mean-pooled) & Rate tokens & 35 \\
Latency profile, $\gamma$ & 30.4\,s (offline rows); 1.04\,s (simultaneous) & Offset $q_0$ / steps $\kappa_h$ / cap $\bar q$ & 2 / 3 per s / 6 \\
Block size $K$ & 5 (one speaker token per 0.4\,s) & Backlog term active from & 0.5\,s; final block $\ge$2 steps \\
Projector hidden size & 1{,}024, output layer zero-initialised & & \\
Time tokens & 251, $\delta_t=80$\,ms, windows $\le$20\,s & \multicolumn{2}{l}{\emph{Playback alignment (\S\ref{sec:method-plem}, \S\ref{sec:method-train})}} \\
Loss weights $\omega_u$ & speaker 2, time 1.5, other 1 & Candidates $|C_i|$ & 8: 4 at base rate, 4 two steps faster \\
LoRA & $r=32$, $\alpha=64$, q/k/v/o, 48 layers & Weights $\eta_\ell$ / $\eta_f$ & 0.5 / 5 (fixed before training) \\
Posterior perturbation & $p=0.2$, $\mathcal{N}(0,0.1^2)$, shift $\le$1 frame & Fidelity $f$ & wav2vec2 CTC, loss clipped at 10 \\
Read interval $\Delta$ & 1\,s & Backlog grid $\delta_h$ / $\bar H$ & 80\,ms / 8\,s ($|\mathcal{H}|=101$) \\
Interpreting lag $\epsilon$ & 1\,s & Budget $b$ & base-rate delay $-$ \num{v10:c:plem:budgetoff}\,s per unit \\
Deadline Thinker rows / steps & \num{v10:c:j:rows} / \num{v10:c:j:s1steps} & $\lambda^\star$ & \num{v10:c:plem:lam:z} (Zh$\rightarrow$En), \num{v10:c:plem:lam:e} (En$\rightarrow$Zh) \\
Learning rate & $10^{-4}\to10^{-5}$, cosine, 100 warm-up & Backlog token & 8 bins of 0.5\,s (capped at 3.5\,s) \\
Batch (global / micro) & 64 (32 per direction) / 8 & $\kappa_q$ / $\kappa_{\mathrm{KL}}$ & \num{v10:c:plem:wrate} / \num{v10:c:plem:wkl} \\
Hardware & 8 H200 GPUs; last step used & Candidate rows kept & $w_i(c)\ge$ \num{v10:c:plem:minq} \\
Decoding & greedy & Test-time samples & 2 per block, highest $f$ kept \\
\bottomrule
\end{tabular}
\caption{Hyperparameters of \sys{}. The per-direction Thinkers of the overlap study (and of the PA Talkers' training
text) follow Table~\ref{tab:hyper2}.}
\label{tab:hyper}
\end{table*}

\subsection{Speaker tokens and the deadline Thinker}
\label{app:train-thinker}
\label{app:model}
\label{app:prompt}

\paragraph{Speaker tokens.} The diarizer's posteriors $\bm{P}\in[0,1]^{F\times S}$ (10\,ms outputs of Nemotron 3
Diarization mean-pooled to 80\,ms frames) are resampled to the acoustic token centres,
$\tilde{\bm{P}}=\mathrm{Interp}(\bm{P};\bar{t}_{1:N})$, zero-padded to $MK$ rows, and every block $\bm{P}_j$ of $K$ rows
becomes one speaker token,
\begin{equation}
\begin{aligned}
\bm{e}_j&=\textstyle\sum_{m=1}^{K}\bm{V}^{(m)}\tilde{\bm{p}}_{(j-1)K+m}+\bm{c},\\
\bm{s}_j&=\bm{V}_2\,\mathrm{GELU}\big(\bm{V}_1\mathrm{GELU}(\bm{e}_j)+\bm{c}_1\big)+\bm{c}_2,
\end{aligned}
\label{eq:conv}
\end{equation}
i.e.\ $g_\psi$ is a strided convolution followed by a two-layer perceptron; $\bm{V}_2,\bm{c}_2$ start at zero, so training starts from
the pretrained Thinker. Per meeting, reference speakers $\varrho$ are mapped to slots by the injective map $\sigma^\star$
that maximises the number of frames in which $\varrho$ speaks and slot $\sigma(\varrho)$ is active ($P>0.5$); an
utterance by $\varrho$ is labelled $\sigma^\star(\varrho)$. Hesitations have no target segment. The instruction asks for
\texttt{<|spk k|><|t a|>}translation\texttt{<|t b|>} per utterance, $k$ the slot, $a,b$ in 80\,ms steps, utterances in
order of start. During training the posteriors are perturbed (Table~\ref{tab:hyper}).

\paragraph{Deadlines.} An utterance's reference translation is split at punctuation into chunks $u_1,\dots,u_J$; chunk
$j$ covers the fraction $\varphi_j=\sum_{j'\le j}|u_{j'}|/\sum_{j'}|u_{j'}|$ of the utterance and ends at the end of
source word $\lceil\varphi_j|\bm{w}|\rceil$ of the time-aligned source words $\bm{w}$, or at $a+\varphi_j(b-a)$ where no
word timings exist (the re-voiced Chinese source). This end time is the $e_j$ of the deadline rule in \S\ref{sec:method-thinker}.

\paragraph{Training.} The deadline Thinker is initialised from a Thinker trained on whole reference windows (the
offline model of Table~\ref{tab:hyper2}, adapters merged). The deadline rows of the two directions are interleaved into
\num{v10:c:j:rows} rows such that every global batch of 64 holds 32 Zh$\rightarrow$En and 32 En$\rightarrow$Zh rows;
there is no checkpoint selection, and the Thinker is the last step. Reads with nothing due are kept with probability
0.5; whole-window (20\%) and single-speaker (15\%) rows are mixed in.

\begin{table}[h]
\centering\small
\setlength{\tabcolsep}{3pt}
\resizebox{\columnwidth}{!}{\begin{tabular}{lcc}
\toprule
 & Offline Thinker & Per-direction Thinker \\
\midrule
Initialisation & Qwen3-Omni Thinker & offline Thinker, merged \\
Directions & both & one model each \\
Diarizer profile & 30.4\,s (offline) & 1.04\,s \\
Rows & 67{,}795 & 128{,}879 / 160{,}957 \\
Steps & 2{,}000 & 3{,}000 \\
Batch (global / micro) & 64 / 8 & 64 / 8 \\
Learning rate & \multicolumn{2}{c}{$10^{-4}\to10^{-5}$, cosine, 50 warm-up} \\
LoRA, loss weights, perturbation & \multicolumn{2}{c}{as in Table~\ref{tab:hyper}} \\
\bottomrule
\end{tabular}}
\caption{The offline Thinker (initialisation of the deadline Thinker) and the per-direction simultaneous Thinkers
(En$\rightarrow$Zh / Zh$\rightarrow$En rows) used in the overlap study and to produce the PA Talkers' training text.}
\label{tab:hyper2}
\end{table}

\subsection{Talker, voice prompts and speaking rate}
\label{app:train-talker}

\paragraph{Voice prompts.} Slot $k$'s solo frames are
\begin{equation}
m_{f,k}=\mathds{1}\big[P_{f,k}>0.5\big]\textstyle\prod_{k'\neq k}\mathds{1}\big[P_{f,k'}<0.5\big].
\label{eq:alone}
\end{equation}
Runs of at least 0.4\,s of solo frames are concatenated into the voice prompt $\bm{v}_k$ (at most 8\,s), using only audio
heard before the slot's first commitment; the Talker encodes it into 32 BiCodec global tokens~\citep{wang2025spark} and
speaks each block as BiCodec semantic tokens, one per 20\,ms.

\paragraph{Speaking rate.} The Talker's rate token takes one of 35 values. B's base Talker uses one fixed value for
every block. In Eq.~\eqref{eq:rate}, the pace $q^{\mathrm{pace}}_i$ is the rate token whose natural duration for the
block's text matches the source time since the previous release; the backlog term is applied once the backlog reaches
0.5\,s, a sentence-final block is raised by at least two steps, and the result is clipped to the token range. Within a
stream, segment $n$ of duration $D_n$ then starts at
\begin{equation}
o_n=\max\big(\zeta_n,\;o_{n-1}+D_{n-1}\big),
\label{eq:layout}
\end{equation}
with $\zeta_n$ its release time; this is the queue of Eq.~\eqref{eq:backlog} before discretisation.

\subsection{Playback alignment and the PA Talker}
\label{app:plem}
\label{app:train-pa}

\begin{figure}[h]
\small
\resizebox{\columnwidth}{!}{\begin{tabular}{@{}l@{}}
\toprule
\textbf{Algorithm 1}: training the PA Talker (one direction) \\
\midrule
1: decode training meetings with a simultaneous Thinker; merge each read's blocks per speaker \\
2: play every stream at the base rate; record $\nu_i$ and the backlog $\tilde{h}_i$ \\
3: for every block, sample $|C_i|$ candidates from the base Talker; compute $D$, $\bar{\ell}$, $f$ \\
4: bisect $\lambda$ until $\mathbb{E}_{p_\lambda}[\Phi_\mathcal{S}]=bW$ (Eq.~\eqref{eq:fb}; Proposition~\ref{prop:pr}) \\
5: run the backward recursion at $\lambda^\star$; weight candidates by $w_i(c)$ (Eq.~\eqref{eq:w}) \\
6: fine-tune the Talker on $\mathcal{J}$ (Eq.~\eqref{eq:patrain}) \\
\textbf{Test}: backlog token of the actual queue; 2 samples per block; keep the higher $f$ \\
\bottomrule
\end{tabular}}
\caption{Training and decoding of the playback layer. No CRF runs at test time.}
\label{alg:pa}
\end{figure}

\paragraph{Candidates and fidelity.} For every block the base Talker samples 8 candidates: 4 at the base rate and 4
at two rate steps faster. $f(c)=-\min(\text{CTC loss per character},10)$ under an independent wav2vec2 CTC
recogniser~\citep{baevski2020wav2vec} of the target language. The forward and backward recursions agree with
brute-force enumeration on small streams.

\paragraph{Objective.} The Talker, conditioned on a backlog token (0.5\,s bins, capped at 3.5\,s), minimises
Eq.~\eqref{eq:patrain} with $\kappa_q=$ \num{v10:c:plem:wrate} and $\kappa_{\mathrm{KL}}=$ \num{v10:c:plem:wkl}; the KL
term covers the semantic codes and the end token, only candidates with $w_i(c)\ge$ \num{v10:c:plem:minq} are kept, and
all Talker parameters are trained.

\paragraph{Data.} Zh$\rightarrow$En: \num{v10:c:plem:meet:z} AMI training meetings, decoded by a separately trained
single-direction simultaneous Thinker (Table~\ref{tab:hyper2}), \num{v10:c:plem:blocks:z} blocks,
\num{v10:c:plem:rows:z} weighted candidate rows, $\lambda^\star=$ \num{v10:c:plem:lam:z}, \num{v10:c:plem:steps:z}
steps. En$\rightarrow$Zh: the same recipe, \num{v10:c:plem:rows:e} rows, $\lambda^\star=$ \num{v10:c:plem:lam:e}; of its
checkpoints at iterations 400, 800 and 1{,}200 a rule on five AMI development meetings (highest tcpASR-BLEU among
checkpoints within 0.1\,s of the last one's speech latency) chose iteration \num{v10:c:plem:it:e}. The budget is the
base-rate path's own delay minus \num{v10:c:plem:budgetoff}\,s per unit.

\paragraph{Ablation rows.} In Table~\ref{tab:ablation}, R1 merges the blocks of one read, R2 speaks them with the base
Talker at the rate of Eq.~\eqref{eq:rate}, R3 replaces the base Talker by the PA Talker with one sample, and R4
(\sys{}) keeps the better of two samples by $f$; R3 $-$ R2 isolates the PA Talker and R4 $-$ R3 the selection.

\section{Inference and Online Processing}
\label{app:online}
\label{app:chunk}
\label{app:causal}

\paragraph{Online windows.} A test meeting is processed whole. A window opens at the first frame where any diarizer
slot is active and closes after 0.6\,s with no active slot or at 20\,s. Slots persist across windows; the decoder
restarts at each window, reading every $\Delta$ seconds as in \S\ref{sec:method-thinker}. All systems of
Table~\ref{tab:main} translate the same windows.

\paragraph{Chunk protocol.} The diarizer's 1.04\,s profile emits frames in blocks of nine; frame $f$ (80\,ms) is known
at $t^{\mathrm{rel}}_f=(9\lfloor f/9\rfloor+13)\cdot 80\,\text{ms}\le t_f+\gamma$, $\gamma=1.04$\,s. A window's opening is
released at $t^{\mathrm{rel}}$ of its first active frame and its end at $t^{\mathrm{rel}}$ of the last frame of the
closing gap; commits before the opening are delayed to it, and those of the last read to the release of the end. All
latencies are scored under this protocol.

\paragraph{Causality.} Simultaneous runs read the posteriors of one run of the diarizer's 1.04\,s profile over the
whole meeting, with every frame later than $t-\gamma$ zeroed at a read at time $t$. For three windows of every test
meeting we truncated the audio at a read time and re-ran the diarizer from scratch: over \num{v6:caus:cuts} cuts the
largest absolute posterior difference was \num{v6:caus:maxabs} (fp16 precision), and every differing $>0.5$ decision
was an fp16 rounding tie within $5\cdot10^{-4}$ of 0.5. The audio encoder encodes the heard prefix at every read.

\section{Playback Alignment: Derivations and Proofs}
\label{app:proofs}

\begin{figure}[t]
\centering
\begin{tikzpicture}[x=0.82cm,y=0.6cm,font=\scriptsize,
  st/.style={circle,draw,inner sep=0pt,minimum size=4.2mm},
  on/.style={st,fill=trained!25},
  blk/.style={rectangle,draw,rounded corners=1pt,inner sep=1.5pt,minimum height=3.2mm}]
% time axis with release times
\draw[->,gray] (-0.3,-3.75) -- (7.4,-3.75) node[right,black]{$t$};
\foreach \x/\l in {0/\nu_i,2.4/\nu_{i+1},4.8/\nu_{i+2}}{\draw[gray] (\x,-3.65)--(\x,-3.85) node[below,black]{$\l$};}
% backlog states (rows = backlog values)
\foreach \c/\x in {i/0,i+1/2.4,i+2/4.8}{\node[gray] at (\x,1.05) {$h_{\c}$};}
\foreach \i/\x in {0/0,1/2.4,2/4.8}{\foreach \j in {0,1,2}{\node[st] (s\i\j) at (\x,-\j) {};}}
\node[left,gray] at (-0.25,0) {0\,s};\node[left,gray] at (-0.25,-1) {0.4\,s};\node[left,gray] at (-0.25,-2) {0.8\,s};
% candidate edges from h_i = 0
\draw[->,thick,trained] (s00) -- node[above,sloped]{$c$: short} (s10);
\draw[->,thick,frozen,dashed] (s00) -- node[below,sloped,pos=0.6]{long} (s12);
\draw[->,thin,gray] (s10) -- (s20); \draw[->,thin,gray] (s10) -- (s21);
\draw[->,thin,gray] (s12) -- (s21); \draw[->,thin,gray] (s12) -- (s22);
% backward message
\draw[->,densely dotted,thick,spkc] (5.6,-2.4) to[bend left=12] node[below,pos=0.45]{backward message $\beta_{i+1}(h)$} (0.6,-2.55);
\node[right,align=left] at (5.15,-0.35) {$\phi_i(h,c)=$\\$\eta_\ell\bar\ell(c)+\eta_f f(c)$\\$-\lambda L_i(h,c)$};
\end{tikzpicture}
\caption{Playback alignment for one speaker stream. States are the backlog $h_i$ (time block $i$ waits for earlier
speech); an edge is a candidate rendering $c$ of block $i$, whose duration sets the next backlog
(Eq.~\eqref{eq:backlog}, discretised). A long candidate delays every later block; backward messages bring that cost into the choice
for block $i$.}
\label{fig:pa}
\end{figure}


Figure~\ref{fig:pa} shows the state space of Eq.~\eqref{eq:crf} for one stream $\mathcal{U}$. Throughout, fix
$\lambda$ and write $\phi_i=\phi^\lambda_i$. A feasible plan is determined by its candidates: given $c_{1:I}$, the
transitions $g_i$ fix every backlog from $h_1=0$. We write $e(\bm{z})=\exp\sum_i\phi_i(z_i)$ for the unnormalised
weight of a plan. The messages of \S\ref{sec:method-plem} are
\begin{equation}
\begin{aligned}
\alpha_i(h)&=\textstyle\sum_{c_{<i}:\,h_i=h}\exp\sum_{j<i}\phi_j(z_j),\\
\beta_i(h)&=\textstyle\sum_{c_{\ge i}\mid h_i=h}\exp\sum_{j\ge i}\phi_j(z_j),
\end{aligned}
\label{eq:msg}
\end{equation}
the total weight of feasible prefixes that leave block $i$ with backlog $h$, and of all completions from there.

\begin{proposition}[Exact inference]
\label{prop:fb}
For the CRF of Eq.~\eqref{eq:crf}, the messages of Eq.~\eqref{eq:msg} satisfy the recursions of Eq.~\eqref{eq:fb};
for every $i$, $Z_\lambda=\sum_h\alpha_i(h)\beta_i(h)$, in particular $Z_\lambda=\beta_1(0)$; and the marginals are given
by Eq.~\eqref{eq:fb}. The cost is $O(I\,|\mathcal{H}|\max_i|C_i|)$.
\end{proposition}

\begin{proof}
\emph{Factorisation.} The backlogs after block $i$ are computed from $h_i$ and $c_i,c_{i+1},\dots$ alone: $c_{<i}$
influences them only through $h_i$. Hence the weight of a plan splits at any block,
\[
e(\bm{z})=\Big(\textstyle\prod_{j<i}e^{\phi_j(h_j,c_j)}\Big)\Big(\textstyle\prod_{j\ge i}e^{\phi_j(h_j,c_j)}\Big),
\]
a prefix factor that ends at $h_i$ times a suffix factor that starts at $h_i$ and depends on the prefix only through
$h_i$.
\emph{Partition function.} Summing $e(\bm{z})$ over all feasible plans and grouping them by the value $h$ of $h_i$, the
prefix sums give $\alpha_i(h)$ and the suffix sums $\beta_i(h)$, so $Z_\lambda=\sum_h\alpha_i(h)\beta_i(h)$. At $i=1$,
$\alpha_1(h)=\mathds{1}[h=0]$ gives $Z_\lambda=\beta_1(0)$.
\emph{Recursions.} A prefix ending at $h_{i+1}=h'$ is a prefix ending at some $h_i=h$ followed by a candidate $c$ with
$g_i(h,c)=h'$; its weight is the old weight times $e^{\phi_i(h,c)}$, which gives $\alpha_{i+1}$. A suffix starting at
$h_i=h$ is a candidate $c$ followed by a suffix starting at $g_i(h,c)$, which gives $\beta_i$.
\emph{Marginals.} The plans with $z_i=(h,c)$ are a prefix ending at $h$, the state $(h,c)$ and a suffix starting at
$g_i(h,c)$, with total weight $\alpha_i(h)e^{\phi_i(h,c)}\beta_{i+1}(g_i(h,c))$; dividing by $Z_\lambda$ gives
Eq.~\eqref{eq:fb}, and $\mathbb{E}_{p_\lambda}[\Phi]=\sum_i\mathbb{E}[L_i(z_i)]$ by linearity.
\emph{Cost.} Each recursion step visits every pair $(h,c)\in\mathcal{H}\times C_i$ once.
\end{proof}

\paragraph{The dual of the latency budget.} The Lagrangian of Eq.~\eqref{eq:pr} is
$\mathrm{KL}(q\Vert p_0)+\lambda(\mathbb{E}_q[\Phi_\mathcal{S}]-bW)$ with multiplier $\lambda\ge0$. Substituting its
minimiser (Proposition~\ref{prop:pr}) gives the scalar dual function
\begin{equation}
\mathcal{D}(\lambda)=-\sum_{\mathcal{U}\in\mathcal{S}}\log\frac{Z_\lambda(\mathcal{U})}{Z_0(\mathcal{U})}-\lambda bW,
\label{eq:dual}
\end{equation}
and, since $\partial_\lambda\log Z_\lambda=-\mathbb{E}_{p_\lambda}[\Phi]$,
\begin{equation}
\mathcal{D}'(\lambda)=\mathbb{E}_{p_\lambda}[\Phi_\mathcal{S}]-bW,\quad
\mathcal{D}''(\lambda)=-\mathrm{Var}_{p_\lambda}[\Phi_\mathcal{S}]\le0.
\label{eq:dualderiv}
\end{equation}

\begin{proposition}[Latency budget]
\label{prop:pr}
Assume some distribution meets the budget of Eq.~\eqref{eq:pr} strictly (it suffices that the minimum-delay plans do).
Then the solution of Eq.~\eqref{eq:pr} is $q^\star=\prod_{\mathcal{U}\in\mathcal{S}}p_{\lambda^\star}(\cdot\mid\mathcal{U})$,
where $\lambda^\star$ maximises the dual of Eq.~\eqref{eq:dual} over $\lambda\ge0$; the dual has the derivatives of
Eq.~\eqref{eq:dualderiv}, and $\lambda^\star=0$ if $\mathbb{E}_{p_0}[\Phi_\mathcal{S}]\le bW$, otherwise
$\mathbb{E}_{p_{\lambda^\star}}[\Phi_\mathcal{S}]=bW$.
\end{proposition}

\begin{proof}
\emph{Inner minimisation.} For $\lambda\ge0$, define the Gibbs distribution
$q_\lambda(\bm{z})=p_0(\bm{z})e^{-\lambda\Phi_\mathcal{S}(\bm{z})}/M_\lambda$ with $M_\lambda=\mathbb{E}_{p_0}[e^{-\lambda\Phi_\mathcal{S}}]$.
Rearranging the Lagrangian,
\[
\begin{aligned}
&\mathrm{KL}(q\Vert p_0)+\lambda\,\mathbb{E}_q[\Phi_\mathcal{S}]\\
&\quad=\textstyle\sum_{\bm{z}}q\log\big(q/(p_0e^{-\lambda\Phi_\mathcal{S}})\big)\\
&\quad=\mathrm{KL}(q\Vert q_\lambda)-\log M_\lambda ,
\end{aligned}
\]
which is minimised, uniquely, by $q=q_\lambda$.
\emph{The minimiser is the CRF.} $p_0$ is a product over streams of Eq.~\eqref{eq:crf} at $\lambda=0$, and
$\Phi_\mathcal{S}$ is a sum of per-block delays, so
$p_0(\bm{z})e^{-\lambda\Phi_\mathcal{S}(\bm{z})}=\prod_{\mathcal{U}}\mathds{1}[\bm{z}_{\mathcal{U}}\in\mathcal{Z}(\mathcal{U})]\exp\big(\sum_i\phi^0_i-\lambda L_i\big)/Z_0(\mathcal{U})$;
since $\phi^0_i-\lambda L_i=\phi^\lambda_i$, normalising gives $q_\lambda=\prod_{\mathcal{U}}p_\lambda(\cdot\mid\mathcal{U})$ and
$M_\lambda=\prod_{\mathcal{U}}Z_\lambda(\mathcal{U})/Z_0(\mathcal{U})$.
\emph{Dual and its derivatives.} Substituting $q_\lambda$ into the Lagrangian gives
$\mathcal{D}(\lambda)=-\log M_\lambda-\lambda bW$, i.e.\ Eq.~\eqref{eq:dual}. For one stream,
$\partial_\lambda\log Z_\lambda=\sum_{\bm{z}}(-\Phi(\bm{z}))e^{\sum_i\phi^\lambda_i}/Z_\lambda=-\mathbb{E}_{p_\lambda}[\Phi]$, and
differentiating once more, $\partial^2_\lambda\log Z_\lambda=\mathbb{E}_{p_\lambda}[\Phi^2]-\mathbb{E}_{p_\lambda}[\Phi]^2=\mathrm{Var}_{p_\lambda}[\Phi]$.
Streams are independent under $q_\lambda$, so expectations and variances add, which gives Eq.~\eqref{eq:dualderiv}.
\emph{Optimality.} Eq.~\eqref{eq:pr} minimises a convex function of $q$ under one linear constraint, and a strictly
feasible point exists, so strong duality holds and $q^\star=q_{\lambda^\star}$ with $\lambda^\star=\arg\max_{\lambda\ge0}\mathcal{D}$.
As $\mathcal{D}''\le0$, $\mathcal{D}'(\lambda)=\mathbb{E}_{p_\lambda}[\Phi_\mathcal{S}]-bW$ is non-increasing: if it is
$\le0$ at $\lambda=0$ the maximum is at $\lambda^\star=0$, otherwise at its root, where the budget holds with equality
(complementary slackness). A non-increasing function has its root found by bisection.
\end{proof}

\begin{proposition}[Posterior-weighted training]
\label{prop:patrain}
Let $\pi$ be a distribution over speech-code sequences, $C$ a set of candidates with distinct code sequences,
$\pi(C)=\sum_{c\in C}\pi(c)>0$ and $\tilde{\pi}(c)=\pi(c)/\pi(C)$. For any distribution $w$ on $C$,
$-\sum_{c\in C}w(c)\log\pi(c)=\mathrm{KL}(w\Vert\tilde{\pi})+\mathrm{H}(w)-\log\pi(C)$.
\end{proposition}

\begin{proof}
Write $\log\pi(c)=\log\tilde{\pi}(c)+\log\pi(C)$ and use $\sum_cw(c)=1$:
$-\sum_cw(c)\log\pi(c)=-\sum_cw(c)\log\tilde{\pi}(c)-\log\pi(C)$, and
$-\sum_cw(c)\log\tilde{\pi}(c)=\mathrm{KL}(w\Vert\tilde{\pi})+\mathrm{H}(w)$.
Since $\mathrm{H}(w)$ does not depend on $\pi$, minimising the weighted cross-entropy over $\pi$ moves $\tilde{\pi}$
towards the CRF posterior $w$ and raises the mass $\pi(C)$ on renderings of the kind the CRF scored.
\end{proof}

\section{Evaluation Protocols}
\label{app:eval}
\label{app:metrics}
\label{app:prereg}

\paragraph{tcpBLEU and tcpASR-BLEU.} Reference units are a speaker's scored utterances chained while less than 1\,s
apart; output slots are matched to speakers per meeting by chrF over the meeting's text; each output segment, on the
meeting clock and widened by a 2\,s collar, goes to the annotated utterance it overlaps most: to that utterance's unit if
it belongs to the segment's matched speaker, otherwise to the matched speaker's own unit it overlaps most, and otherwise
it is an insertion. The score is corpus BLEU over (unit hypothesis, unit reference) pairs plus (insertion, empty) pairs:
tcpBLEU on the text output, tcpASR-BLEU on the ASR transcripts of the speech.

\paragraph{cpASR-BLEU.} Per meeting, the ASR transcripts of each output speaker label are concatenated in time order,
and so are the scored reference utterances of each reference speaker (the units of tcpBLEU); a transcript segment is
left out, as in tcpBLEU, when the annotated utterance it overlaps most within the collar is not scored (a hesitation,
or speech outside the reference windows). Matching and scoring follow \S\ref{sec:method-metrics}.

\paragraph{Recognisers and BLEU.} Speech is transcribed by Whisper-large-v3 (English) or Paraformer (Chinese). BLEU is
corpus-level sacreBLEU~\citep{post2018call} (13a for English, zh for Chinese). Our systems are synthesised with two TTS
runs, pooled per meeting as sufficient statistics; external systems have one.

\paragraph{Speech latency.} For a reference unit with source duration $|X|$ and $|Y|$ heard target words (words or
Chinese characters) with delays $d_1\le\dots\le d_{|Y|}$ from the unit's start,
$\mathrm{LAAL}=\frac{1}{k^\ast}\sum_{k=1}^{k^\ast}\big(d_k-(k-1)\frac{|X|}{\max(|Y|,|Y^\ast|)}\big)$, with $|Y^\ast|$ the
reference length and $k^\ast$ the first $k$ with $d_k\ge|X|$~\citep{papi2022over}. Block $i$ of a stream starts
playing at $\nu_i+h_i$. LAAL (LAAL$_\text{block}$ in the appendix tables) gives every word of the block
$d_k=\nu_i+h_i$ relative to the unit's start; LAAL$_\text{spread}$ spreads its $W_i$ words evenly over the block's
duration, as Eq.~\eqref{eq:lat}. Speech start is the first heard word of a unit; EndOffset runs from the unit's end to
the end of its last heard word.

\paragraph{Statistics and analysis plans.} Paired tests are over meetings: exact two-sided sign-flip permutation tests
over all $2^M$ flips and 95\% meeting-bootstrap intervals (10{,}000 resamples). Each study followed an internal analysis
plan written before its test decoding and scoring, with every change recorded as a dated amendment before the results
it touches. The plans fix the meeting lists (AMI test EN2002a--d, ES2004a--d, IS1009a--d, TS3003a--d; all 20 AliMeeting
test meetings; the eight AliMeeting evaluation meetings), the systems, the confirmatory contrasts with their Holm
families (\sys{} $-$ B over twelve BLEU tests; \sys{} $-$ each system of Table~\ref{tab:main} per BLEU score, fixed
before the added systems were decoded) and the development rules for choosing checkpoints. The plans are timestamped
only in a private repository, so reviewers cannot verify them; we call their tests \emph{confirmatory} but do not rely
on them as public pre-registrations.

\section{Baseline Systems}
\label{app:baselines}

\paragraph{Causal relabelling.} Systems without meeting-wide speaker labels (SeamlessStreaming, X-Translator, the
cascade) are attributed per output segment to the diarizer slot most active over the segment, using only frames
released when the segment is committed.

\subsection{SeamlessStreaming}
\label{app:e7}
SeamlessStreaming~\citep{seamless2023} (official agent and default settings, inference only) reads each of our online
test windows; each output segment is attributed by causal relabelling, and text, speech and latency are scored with our
tools under the chunk protocol. Nothing was trained or tuned, and each test set was decoded once. It is a zero-shot
reference for a public single-speaker SimulS2ST system on overlapping meetings, not a like-for-like baseline
(Table~\ref{tab:e7}).
\begin{table}[t]
\centering\small
\setlength{\tabcolsep}{2.5pt}
\resizebox{\columnwidth}{!}{\begin{tabular}{llrrrrr}
\toprule
Set & System & tcpBLEU & ov & tcpASR-BLEU & LAAL$_\text{spread}$ & EndOff. \\
\midrule
AMI Zh$\rightarrow$En & \sys{} & 25.02 & 19.97 & 20.35 & 3.61 & 3.13 \\
 & SeamlessStreaming & 8.69 & 3.07 & 8.48 & 3.29 & 2.53 \\
\midrule
AliMeeting & \sys{} & 21.43 & 16.91 & 18.08 & 3.65 & 3.23 \\
 & SeamlessStreaming & 10.41 & 4.42 & 9.42 & 3.96 & 2.72 \\
\midrule
AMI En$\rightarrow$Zh & \sys{} & 41.15 & 36.16 & 32.23 & 2.95 & 2.91 \\
 & SeamlessStreaming & 13.28 & 5.86 & 11.31 & 3.19 & 3.04 \\
\bottomrule
\end{tabular}}
\caption{External zero-shot single-speaker S2ST reference on our online test windows (no speaker attribution of its own: causal diarizer relabelling; one deterministic run; chunk protocol, seconds). LAAL compares systems of similar quality only, so its latency does not place it on our trade-off.}
\label{tab:e7}
\end{table}


\ifXT
\subsection{X-Translator}
\label{app:xt}
X-Translator~\citep{zhao2026xtranslator} (released code, weights and configuration: Qwen3-ASR-1.7B for simultaneous
recognition, LMT-60-8B translation, X-Voice TTS cloning a voice prompt per tracked speaker) runs each online window as
one fresh session with the per-session logic of its server. A translated segment counts as written when the audio that
produced it has been received (computation is not counted), is queued for playback like ours, and is attributed by
causal relabelling; its own speaker tracker only chooses the cloned voice. Its client skips 0.16\,s chunks with RMS
below 0.01, a gate meant for browser input with automatic gain control. Our AMI audio is quieter (median loudest-chunk
RMS 0.016--0.018, against 0.066 on AliMeeting), and the gate never opens in 32\% / 28\% of the AMI Zh$\rightarrow$En /
En$\rightarrow$Zh test windows. Before any X-Translator score was seen we therefore fixed one constant gain per AMI set
on its dev windows, $g=0.1/m_\text{dev}$ ($\times$1.81 / $\times$2.16, $m_\text{dev}$ the dev median of the
loudest-chunk RMS), which leaves 8\% / 10\% of the test windows gated; AliMeeting (1.4\% gated) is used as is. Without
the gain X-Translator reaches \num{v14:ami:xtraw:cpASR} / \num{v14:en:xtraw:cpASR} cpASR-BLEU and
\num{v14:ami:xtraw:tcpASR} / \num{v14:en:xtraw:tcpASR} tcpASR-BLEU on the two AMI sets. Each configuration was decoded
once.
\fi

\ifCascade
\subsection{Same-data cascade}
\label{app:cascade}
The cascade of Table~\ref{tab:main} asks what translating speaker-blind and attributing afterwards gives with our
training data, detector windows and Talker. Its simultaneous model is the speaker-blind per-direction Thinker of
Table~\ref{tab:overlap} (trained like the per-direction Thinkers of Table~\ref{tab:hyper2} with the speaker tokens
zeroed), decoded online with B's settings; each segment is relabelled causally and spoken by the base Talker with B's
settings (two TTS runs). Without relabelling, its own guessed labels give text tcpBLEU of only 10.9 / 24.4 / 11.4 (AMI
Zh$\rightarrow$En / En$\rightarrow$Zh / AliMeeting), so its attribution comes from the diarizer. After relabelling its
text is close to B's on single-speaker units but lower on overlapped units (15.7 vs.\ 20.0 on AMI Zh$\rightarrow$En,
30.9 vs.\ 36.2 on En$\rightarrow$Zh), the gap that speaker tokens close in Table~\ref{tab:overlap}.
\fi

\section{Supplementary Results}
\label{app:more}
\label{app:sig10}
\label{app:postreview}

\subsection{Significance of Table~\ref{tab:main}}
Table~\ref{tab:syssig} lists the paired tests behind Table~\ref{tab:main}, and Figure~\ref{fig:tradeoff} plots its
speech quality against LAAL together with the ablation path from B to \sys{}. All twelve BLEU contrasts of each score have
$p_\text{Holm}<0.001$. In LAAL, \sys{} is heard earlier than the cascade and B on every set and earlier than both public systems on
AMI En$\rightarrow$Zh (SeamlessStreaming: not significant); it is heard later than X-Translator on AMI
Zh$\rightarrow$En and AliMeeting and than SeamlessStreaming on AMI Zh$\rightarrow$En.
\begin{table*}[t]
\centering\small
\setlength{\tabcolsep}{3pt}
\begin{tabular}{llrrrrrr}
\toprule
Set & \sys{} $-$ & cpASR-BLEU [95\%] & $p_\text{Holm}$ & tcpASR-BLEU [95\%] & $p_\text{Holm}$ & LAAL (s) [95\%] & $p$ \\
\midrule
AMI Zh$\rightarrow$En & SeamlessStreaming & +14.75 [13.43, 15.92] & $<$0.001 & +11.87 [11.29, 12.46] & $<$0.001 & +0.30 [0.16, 0.46] & 0.003 \\
 & X-Translator & +14.79 [12.74, 16.76] & $<$0.001 & +11.71 [11.00, 12.38] & $<$0.001 & +0.34 [0.16, 0.54] & 0.004 \\
 & Cascade & +1.97 [1.60, 3.16] & $<$0.001 & +1.98 [1.57, 3.08] & $<$0.001 & $-$0.48 [$-$0.56, $-$0.41] & $<$0.001 \\
 & B & +2.12 [1.25, 3.08] & $<$0.001 & +2.28 [1.58, 3.04] & $<$0.001 & $-$0.09 [$-$0.10, $-$0.08] & $<$0.001 \\
\midrule
AMI En$\rightarrow$Zh & SeamlessStreaming & +22.69 [21.49, 23.81] & $<$0.001 & +20.91 [19.87, 21.86] & $<$0.001 & $-$0.07 [$-$0.14, 0.01] & 0.091 \\
 & X-Translator & +13.18 [10.90, 15.56] & $<$0.001 & +12.27 [10.28, 14.35] & $<$0.001 & $-$0.31 [$-$0.41, $-$0.22] & $<$0.001 \\
 & Cascade & +1.41 [0.93, 1.85] & $<$0.001 & +1.26 [0.79, 1.66] & $<$0.001 & $-$0.37 [$-$0.46, $-$0.31] & $<$0.001 \\
 & B & +0.66 [0.48, 0.84] & $<$0.001 & +0.58 [0.43, 0.74] & $<$0.001 & $-$0.11 [$-$0.13, $-$0.09] & $<$0.001 \\
\midrule
AliMeeting Zh$\rightarrow$En & SeamlessStreaming & +10.45 [8.79, 12.14] & $<$0.001 & +8.66 [7.52, 9.77] & $<$0.001 & $-$0.29 [$-$0.48, $-$0.11] & 0.005 \\
 & X-Translator & +4.58 [3.09, 5.99] & $<$0.001 & +4.31 [3.22, 5.29] & $<$0.001 & +0.18 [0.07, 0.30] & 0.009 \\
 & Cascade & +2.23 [1.93, 2.99] & $<$0.001 & +1.92 [1.64, 2.54] & $<$0.001 & $-$0.60 [$-$0.71, $-$0.50] & $<$0.001 \\
 & B & +1.44 [1.07, 2.60] & $<$0.001 & +1.35 [1.09, 2.29] & $<$0.001 & $-$0.33 [$-$0.40, $-$0.28] & $<$0.001 \\
\bottomrule
\end{tabular}
\caption{Paired tests of Table~\ref{tab:main}: \sys{} minus each system over meetings (exact sign-flip test, 95\% meeting-bootstrap interval; BLEU contrasts Holm-corrected per score as fixed in the analysis plan; LAAL uncorrected). A positive LAAL difference means \sys{} is heard later.}
\label{tab:syssig}
\end{table*}

\begin{figure*}[t]
\centering
\begin{tikzpicture}
\begin{groupplot}[group style={group size=3 by 1, horizontal sep=1.1cm}, width=0.36\textwidth, height=3.2cm, xlabel={LAAL (s) $\downarrow$}, title style={font=\small, yshift=-1mm}, tick label style={font=\scriptsize}, label style={font=\small}, grid=major, grid style={gray!20}, legend style={font=\scriptsize, draw=none, column sep=4pt}]
\nextgroupplot[title={AMI Zh$\rightarrow$En}, ylabel={tcpASR-BLEU $\uparrow$}, legend to name=leg:tradeoff, legend columns=6]
\addplot[only marks, mark=triangle*, mark size=2.6pt, color=gray!70!black] coordinates {(2.73,8.48)};\addlegendentry{SeamlessStreaming}
\addplot[only marks, mark=diamond*, mark size=2.6pt, color=frozen] coordinates {(2.69,8.64)};\addlegendentry{X-Translator}
\addplot[only marks, mark=square*, mark size=2.6pt, color=spkc] coordinates {(3.51,18.37)};\addlegendentry{Cascade}
\addplot[only marks, mark=*, mark size=2.6pt, color=black] coordinates {(3.12,18.07)};\addlegendentry{B}
\addplot[only marks, mark=pentagon*, mark size=2.6pt, color=trained] coordinates {(3.03,20.35)};\addlegendentry{\protect\sys{}}
\addplot[thin, color=trained, mark=o, mark size=1.2pt, -{Stealth[length=1.6mm]}] coordinates {(3.12,18.07) (3.09,18.32) (3.04,19.88) (3.03,20.17) (3.03,20.35)};\addlegendentry{B, R1--R3, \protect\sys{} (ablation path)}
\nextgroupplot[title={AMI En$\rightarrow$Zh}]
\addplot[only marks, mark=triangle*, mark size=2.6pt, color=gray!70!black] coordinates {(2.50,11.31)};
\addplot[only marks, mark=diamond*, mark size=2.6pt, color=frozen] coordinates {(2.74,19.96)};
\addplot[only marks, mark=square*, mark size=2.6pt, color=spkc] coordinates {(2.80,30.97)};
\addplot[only marks, mark=*, mark size=2.6pt, color=black] coordinates {(2.54,31.65)};
\addplot[only marks, mark=pentagon*, mark size=2.6pt, color=trained] coordinates {(2.43,32.23)};
\addplot[thin, color=trained, mark=o, mark size=1.2pt, -{Stealth[length=1.6mm]}] coordinates {(2.54,31.65) (2.50,31.67) (2.45,31.91) (2.43,31.93) (2.43,32.23)};
\nextgroupplot[title={AliMeeting Zh$\rightarrow$En (out-of-domain)}]
\addplot[only marks, mark=triangle*, mark size=2.6pt, color=gray!70!black] coordinates {(3.24,9.42)};
\addplot[only marks, mark=diamond*, mark size=2.6pt, color=frozen] coordinates {(2.77,13.77)};
\addplot[only marks, mark=square*, mark size=2.6pt, color=spkc] coordinates {(3.55,16.16)};
\addplot[only marks, mark=*, mark size=2.6pt, color=black] coordinates {(3.28,16.73)};
\addplot[only marks, mark=pentagon*, mark size=2.6pt, color=trained] coordinates {(2.95,18.08)};
\addplot[thin, color=trained, mark=o, mark size=1.2pt, -{Stealth[length=1.6mm]}] coordinates {(3.28,16.73) (3.16,16.62) (3.00,17.49) (2.95,17.82) (2.95,18.08)};
\end{groupplot}
\end{tikzpicture}\\[-1mm]
\pgfplotslegendfromname{leg:tradeoff}
\caption{Speech quality against speech LAAL for the systems of Table~\ref{tab:main} (upper left is better). The line traces B through the playback-layer components of Table~\ref{tab:ablation} to \sys{} (points: mean of the two TTS runs).}
\label{fig:tradeoff}
\end{figure*}


\subsection{\sys{} against B}
Table~\ref{tab:sig} lists the confirmatory tcpASR-BLEU tests of \sys{} against B, whose text is identical.
Table~\ref{tab:speech} gives both speech-LAAL definitions, speech start, EndOffset, speaking rate,
UTMOS~\citep{saeki2022utmos}, speaker similarity (WavLM, against the Talker's voice prompt) and pauses, with paired
latency differences and their intervals. Every latency definition falls; the speaking rate rises by
\num{v11:pr:test:ami:e2-b:ratepct} / \num{v11:pr:test:ali:e2-b:ratepct} / \num{v11:pr:test:en:e2-b:ratepct}\%
(AMI Zh$\rightarrow$En / AliMeeting / En$\rightarrow$Zh); UTMOS, paired per stream, rises by
\num{v10:pr:test:ami:e2-b:utmos:diff} / \num{v10:pr:test:ali:e2-b:utmos:diff} / \num{v10:pr:test:en:e2-b:utmos:diff};
speaker similarity is unchanged; pauses become two to three times as frequent. Table~\ref{tab:decompfull} adds the
development sets to the delay decomposition of Table~\ref{tab:decomp}.
\begin{table}[t]
\centering\small
\setlength{\tabcolsep}{2.5pt}
\resizebox{\columnwidth}{!}{\begin{tabular}{llrrrr}
\toprule
Set & Comparison & $\Delta$ & 95\% interval & $p$ & $p_\text{Holm}$ \\
\midrule
AMI Zh$\rightarrow$En & tcpASR-BLEU: \sys{} $-$ B & +2.28 & [1.57, 3.04] & $<$0.001 & $<$0.001 \\
\midrule
AliMeeting & tcpASR-BLEU: \sys{} $-$ B & +1.35 & [1.09, 2.28] & $<$0.001 & $<$0.001 \\
\midrule
AMI En$\rightarrow$Zh & tcpASR-BLEU: \sys{} $-$ B & +0.58 & [0.43, 0.74] & $<$0.001 & $<$0.001 \\
\bottomrule
\end{tabular}}
\caption{Paired tests over meetings on the test sets (16 AMI, 20 AliMeeting meetings): exact two-sided sign-flip permutation $p$ over all $2^M$ flips, 95\% meeting-bootstrap interval (10{,}000 resamples), Holm over the twelve BLEU tests of the analysis plan. tcpASR-BLEU pools the two TTS seeds per meeting. B: deadline baseline; \sys{}: its text with rl + PA Talker. The text of \sys{} equals B's by construction. The intervals do not include decoder run-to-run noise; training-seed variance is not covered.}
\label{tab:sig}
\end{table}

\begin{table*}[t]
\centering\small
\setlength{\tabcolsep}{3.5pt}
\begin{tabular}{llrrrrrrrrr}
\toprule
 & & \multicolumn{4}{c}{speech latency (s)} & & \multicolumn{2}{c}{voice} & \multicolumn{2}{c}{pauses / min} \\
\cmidrule(lr){3-6}\cmidrule(lr){8-9}\cmidrule(lr){10-11}
Set & Sys. & LAAL$_\text{block}$ & LAAL$_\text{spread}$ & start & EndOffset & rate & UTMOS & SIM & $\ge$0.5\,s & mid-sent. \\
\midrule
AMI Zh$\rightarrow$En & B & 3.12 & 4.01 & 2.99 & 4.08 & 2.40 & 1.80 & 0.827 & 3.3 & 0.8 \\
 & \sys{} & 3.03 & 3.61 & 2.96 & 3.13 & 3.49 & 1.94 & 0.827 & 8.3 & 2.8 \\
\midrule
AliMeeting & B & 3.28 & 4.28 & 3.05 & 4.57 & 2.40 & 3.48 & 0.873 & 4.4 & 0.9 \\
 & \sys{} & 2.95 & 3.65 & 2.92 & 3.23 & 3.56 & 3.72 & 0.873 & 11.8 & 3.9 \\
\midrule
AMI En$\rightarrow$Zh & B & 2.54 & 3.21 & 2.43 & 3.48 & 3.90 & 3.25 & 0.833 & 4.6 & -- \\
 & \sys{} & 2.43 & 2.95 & 2.39 & 2.91 & 5.04 & 3.40 & 0.831 & 9.4 & -- \\
\bottomrule
\end{tabular}
\caption{Speech of B and \sys{} on the test sets. Latency: chunk protocol, mean of two TTS seeds; block-start LAAL charges every word of a block when the block starts playing, LAAL$_\text{spread}$ spreads them over its playback; start: first heard word of a unit. Rate: English words (Zh$\rightarrow$En) or characters (En$\rightarrow$Zh) per second of synthesised speech. UTMOS and SIM (WavLM speaker similarity to the Talker's voice prompt, not to the speaker's own recording): one TTS seed, all streams of at least 0.5\,s. Pauses: silences of at least 0.5\,s per minute of speech; mid-sentence: after a block that does not end a sentence (English output only).}
\label{tab:speech}
\end{table*}

\begin{table*}[t]
\centering\small
\setlength{\tabcolsep}{3pt}
\resizebox{\textwidth}{!}{\begin{tabular}{llllll}
\toprule
Set & Comparison & LAAL$_\text{block}$ & LAAL$_\text{spread}$ & start & EndOffset \\
\midrule
AMI Zh$\rightarrow$En & \sys{} $-$ B & $-$0.09 [$-$0.10, $-$0.08] & $-$0.40 [$-$0.42, $-$0.38] & $-$0.03 [$-$0.04, $-$0.02] & $-$0.96 [$-$1.01, $-$0.91] \\
\midrule
AliMeeting & \sys{} $-$ B & $-$0.33 [$-$0.40, $-$0.28] & $-$0.63 [$-$0.69, $-$0.57] & $-$0.13 [$-$0.17, $-$0.10] & $-$1.34 [$-$1.45, $-$1.23] \\
\midrule
AMI En$\rightarrow$Zh & \sys{} $-$ B & $-$0.11 [$-$0.13, $-$0.09] & $-$0.26 [$-$0.28, $-$0.24] & $-$0.04 [$-$0.06, $-$0.03] & $-$0.58 [$-$0.62, $-$0.54] \\
\bottomrule
\end{tabular}}
\caption{Paired speech-latency differences on the test sets (seconds): difference and 95\% meeting-bootstrap interval (10{,}000 resamples, two TTS seeds pooled per meeting). Exact sign-flip $p$ values are given in the text where used; no multiple-test correction.}
\label{tab:speech2}
\end{table*}

\begin{table}[t]
\centering\small
\setlength{\tabcolsep}{2.5pt}
\resizebox{\columnwidth}{!}{\begin{tabular}{llrrrrrrr}
\toprule
Set & Sys. & start & model & detector & queue & last block & EndOff. & load \\
\midrule
dev Zh$\rightarrow$En & B & 2.95 & 0.74 & 0.61 & 0.64 & 2.16 & 4.14 & 1.41 \\
 & \sys{} & 2.90 & 0.75 & 0.61 & 0.12 & 1.62 & 3.10 & 0.94 \\
\midrule
dev En$\rightarrow$Zh & B & 2.36 & 0.76 & 0.55 & 0.59 & 1.54 & 3.45 & 1.33 \\
 & \sys{} & 2.32 & 0.78 & 0.55 & 0.14 & 1.37 & 2.83 & 1.00 \\
\midrule
AMI Zh$\rightarrow$En & B & 2.99 & 0.67 & 0.66 & 0.65 & 2.10 & 4.08 & 1.37 \\
 & \sys{} & 2.96 & 0.68 & 0.66 & 0.14 & 1.64 & 3.13 & 0.94 \\
\midrule
AliMeeting & B & 3.05 & 0.88 & 0.40 & 1.33 & 1.96 & 4.57 & 1.56 \\
 & \sys{} & 2.92 & 0.95 & 0.40 & 0.33 & 1.55 & 3.23 & 1.05 \\
\midrule
AMI En$\rightarrow$Zh & B & 2.43 & 0.69 & 0.64 & 0.54 & 1.61 & 3.48 & 1.28 \\
 & \sys{} & 2.39 & 0.71 & 0.64 & 0.14 & 1.42 & 2.91 & 0.99 \\
\bottomrule
\end{tabular}}
\caption{Development and test sets: where the speech delay of a source unit comes from (chunk protocol, mean over units and two TTS seeds, seconds). EndOffset = model lag (last commit after the unit's end) + detector wait (online window release) + queue wait (earlier speech of the same stream) + duration of the last block. start: speech start offset from the unit's start. load: synthesised speech per second of source speech. \sys{}: the deadline Thinker's text with read-level merging and the PA Talker.}
\label{tab:decompfull}
\end{table}


\ifAliEval
\subsection{Locked test on unused AliMeeting meetings}
\label{app:alieval}
The eight AliMeeting evaluation meetings (\num{v11:ae:hours}\,h) were never used for training, development or any
earlier test. Windows, diarizer runs and scoring follow the AliMeeting test set, except that no window is filtered by
its reference. The analysis plan fixed the systems and three confirmatory hypotheses (\sys{} $-$ B on tcpASR-BLEU, LAAL
and EndOffset, Holm over three) before decoding; each system was decoded once (Table~\ref{tab:alieval}).
\begin{table*}[t]
\centering\small
\setlength{\tabcolsep}{3.5pt}
\begin{tabular}{lrrrrrrrrrr}
\toprule
 & text & tcpASR- & text & \multicolumn{4}{c}{speech latency (s)} & rate & pauses & \\
\cmidrule(lr){5-8}
System & tcpBLEU & BLEU & LAAL & LAAL$_\text{block}$ & LAAL$_\text{spread}$ & start & EndOff. & (w/s) & /min & UTMOS \\
\midrule
B & 14.58 & 10.66 & 2.87 & 3.39 & 4.50 & 3.23 & 4.79 & 2.36 & 4.3 & 3.50 \\
\sys{} & 14.58 & 12.89 & 2.87 & 3.02 & 3.81 & 3.09 & 3.35 & 3.50 & 11.1 & 3.77 \\
\bottomrule
\end{tabular}
\caption{Locked test on the eight unused AliMeeting evaluation meetings (4.2\,h, Zh$\rightarrow$En, DiariST human references), each system decoded once; tcpASR-BLEU and speech metrics averaged over two TTS seeds (UTMOS: seed 1), latencies under the chunk protocol without computation. Confirmatory tests of \sys{} $-$ B (Holm over three; exact sign-flip over 8 meetings, minimum $p$ = 2/256; 95\% meeting-bootstrap interval): tcpASR-BLEU +2.23 [+0.75, +3.37], $p_\text{Holm}$ 0.023; LAAL$_\text{block}$ $-$0.37 [$-$0.41, $-$0.33], $p_\text{Holm}$ 0.023; EndOffset $-$1.43 [$-$1.51, $-$1.32], $p_\text{Holm}$ 0.023.}
\label{tab:alieval}
\end{table*}

\fi

\ifMetricsV
\subsection{A second recogniser and COMET}
\label{app:metricsv11}
Candidate selection uses a CTC recogniser (wav2vec~2.0), and tcpASR-BLEU uses Whisper-large-v3 or Paraformer, so part
of the gain could be speech that suits recognisers rather than listeners. We rescored the existing test outputs with
definitions fixed before computing: the speech of both TTS runs transcribed by a second recogniser of a different
family (English: Parakeet-TDT-0.6B-v2; Chinese: Whisper-large-v3, whose repetition loops and video-outro hallucinations
are removed by one rule for all systems), and COMET~\citep{rei2022comet} per reference unit (Table~\ref{tab:metricsv11}).
The En$\rightarrow$Zh second-recogniser gain is larger than with Paraformer because Whisper fails more often on B's
audio (loops or hallucinations in 16\% of B's utterances against 6\% of \sys{}'s), so only its sign is informative.
\begin{table*}[t]
\centering\small
\setlength{\tabcolsep}{4pt}
\begin{tabular}{lrrrrrrrrr}
\toprule
 & \multicolumn{2}{c}{tcpASR-BLEU} & \multicolumn{3}{c}{tcpASR-BLEU, second recogniser} & \multicolumn{3}{c}{COMET of speech} & COMET \\
\cmidrule(lr){2-3}\cmidrule(lr){4-6}\cmidrule(lr){7-9}
Test set & B & \sys{} & B & \sys{} & $\Delta$ & B & \sys{} & $\Delta$ & of text \\
\midrule
AMI Zh$\rightarrow$En & 18.07 & 20.35 & 19.30 & 20.61 & +1.32$^{*}$ & 73.86 & 76.44 & +2.58$^{*}$ & 79.42 \\
AMI En$\rightarrow$Zh & 31.65 & 32.23 & 28.13 & 31.87 & +3.74$^{*}$ & 73.16 & 73.50 & +0.34$^{*}$ & 78.81 \\
AliMeeting & 16.73 & 18.08 & 17.60 & 18.54 & +0.93$^{*}$ & 64.97 & 66.44 & +1.47$^{*}$ & 71.49 \\
\bottomrule
\end{tabular}
\caption{The speech-layer gain under a second recogniser (English: Parakeet-TDT-0.6B-v2, Chinese: Whisper-large-v3; neither is the CTC recogniser used for selection; repetition loops and video-outro hallucinations of the second recogniser removed by one fixed rule for all systems) and under COMET (wmt22-comet-da, $\times100$, per reference unit, empty hypotheses scored). tcpASR-BLEU: the paper's recognisers (Whisper-large-v3 / Paraformer). Both TTS seeds pooled; B and \sys{} share their text. $\Delta$: \sys{} $-$ B, meeting sign-flip test, Holm over the test family: $^{*}$ $p_\text{Holm}<0.05$, $^{\circ}$ not significant.}
\label{tab:metricsv11}
\end{table*}

\fi

\subsection{Text: DiariST-AliMeeting with its official metric}
\label{app:diarist}
DiariST~\citep{yang2024diarist} reports speaker-attributed text translation on AliMeeting with its own split into 3--6
minute mini-sessions and its own metrics. We cut our whole-meeting outputs into the official mini-sessions (each
segment by its start time) and scored them with the released evaluation script, unmodified (Table~\ref{tab:diarist}).
The audio condition (IHM-MIX) and the human references are those of our AliMeeting test; the published systems are
trained on AliMeeting and decode each mini-session separately, ours are not trained on it and decode whole meetings
online. \sys{}'s text leads on test and is level on dev; its speech, scored on ASR transcripts, stays above the
published text systems and X-Translator's speech on test.
\begin{table}[t]
\centering\small
\setlength{\tabcolsep}{2.5pt}
\resizebox{\columnwidth}{!}{\begin{tabular}{lllrrrr}
\toprule
 & & & \multicolumn{2}{c}{dev} & \multicolumn{2}{c}{test} \\
\cmidrule(lr){4-5}\cmidrule(lr){6-7}
System & output & mode & SAg & SAt & SAg & SAt \\
\midrule
Whisper-small, SD$\rightarrow$ST$^\dagger$ & text & offline & 10.71 & 9.21 & 10.40 & 9.31 \\
Whisper-small, ST$\rightarrow$SD$^\dagger$ & text & offline & 14.48 & 11.96 & 13.68 & 12.10 \\
DiariST (t-SOT + t-vector)$^\dagger$ & text & simultaneous & 16.05 & 13.91 & 16.62 & 14.76 \\
\midrule
\sys{} Thinker (= B) & text & simultaneous & 15.26 & 13.74 & 22.38 & 21.06 \\
\sys{} & speech & simultaneous & 14.54 & 12.78 & 20.25 & 18.61 \\
B & speech & simultaneous & 12.50 & 10.95 & 19.02 & 17.38 \\
SeamlessStreaming & speech & simultaneous & -- & -- & 9.28 & 7.97 \\
X-Translator & speech & simultaneous & -- & -- & 16.04 & 13.66 \\
\bottomrule
\end{tabular}}
\caption{DiariST-AliMeeting (IHM-MIX, human English references) with its official evaluation: speaker-agnostic (SAg) and speaker-attributed (SAt) BLEU over 76 dev / 195 test mini-sessions of 3--6 min. $^\dagger$Published by \citet{yang2024diarist}; these systems are trained on AliMeeting, ours are not. Ours decode whole meetings simultaneously and are cut into the mini-sessions afterwards; speech rows score the ASR transcripts (two TTS seeds averaged for B and \sys{}).}
\label{tab:diarist}
\end{table}


\subsection{Overlap study}
\label{app:v8}
\label{sec:eval-overlap}
A tcpBLEU unit is \emph{overlapped} when another speaker talks during at least 20\% of it. \emph{Concurrent-speaker
recall} (CR) takes every pair of utterances of two speakers that overlap by at least 0.5\,s and asks whether both
received text attributed to their own speaker within 1\,s. The study's analysis plan (internal) was fixed before any of
its test decodes; it fixes the systems, the overlap metrics and their thresholds, the causal relabelling rule and twelve
confirmatory tests (Table~\ref{tab:reg8}). Table~\ref{tab:ovmain8} adds offline decoding on reference windows to the
simultaneous rows of Table~\ref{tab:overlap} (there: ``with speaker tokens'' = \sys{}-Simul, ``speaker-blind +
relabel'' = speaker-blind Simul + causal relabel). Relabelling the offline Thinker's own output with the diarizer,
keeping its text, separates labels from translation: on overlapped units its own labels score
\num{v8:ov:en2zh:O2:ov} against \num{v8:ov:en2zh:O2-relabel:ov} with diarizer labels (En$\rightarrow$Zh), while with
diarizer labels on both sides its text and the speaker-blind model's are close
(\num{v8:ov:en2zh:O2-relabel:ov} vs.\ \num{v8:ov:en2zh:O2-nosd-relabel:ov}).

\begin{table*}[t]
\centering\small
\setlength{\tabcolsep}{4pt}
\begin{tabular}{lrrrrrrrrr}
\toprule
 & \multicolumn{3}{c}{En$\rightarrow$Zh} & \multicolumn{3}{c}{Zh$\rightarrow$En} & \multicolumn{3}{c}{AliMeeting} \\
\cmidrule(lr){2-4}\cmidrule(lr){5-7}\cmidrule(lr){8-10}
System & all & overlap & CR & all & overlap & CR & all & overlap & CR \\
\midrule
\multicolumn{10}{l}{\emph{Reference windows, offline}} \\
Cascade, same data & 36.13 & 26.21 & 89.3 & 21.96 & 14.36 & 87.4 & 15.70 & 9.79 & 93.3 \\
Speaker-blind + relabel & 41.39 & 30.98 & 51.8 & 24.36 & 16.90 & 50.6 & 16.27 & 12.04 & 58.0 \\
\sys{} (offline) & 42.06 & 34.56 & 76.3 & 25.43 & 19.72 & 79.0 & 15.19 & 12.10 & 79.2 \\
\multicolumn{10}{l}{\emph{Whole meetings, online, 1\,s reads}} \\
Speaker-blind Simul + causal relabel & 35.22 & 27.17 & 65.5 & 19.33 & 9.06 & 54.4 & 12.67 & 8.55 & 63.0 \\
\sys{}-Simul & 40.09 & 35.71 & 82.3 & 24.53 & 17.73 & 80.2 & 11.47 & 9.18 & 80.2 \\
\bottomrule
\end{tabular}
\caption{Overlap-aware evaluation (\S\ref{sec:eval-overlap}): tcpBLEU over all units and over overlapped units, and
concurrent-speaker recall (CR, \%). Relabelling attaches one diarizer speaker to each translated segment. AliMeeting references: machine translations of the human transcripts (Appendix~\ref{app:data}).}
\label{tab:ovmain8}
\end{table*}

\begin{table*}[t]
\centering\small
\setlength{\tabcolsep}{4pt}
\begin{tabular}{lllrrrl}
\toprule
Test & Set & Comparison (whole meetings, online) & $\Delta$ & 95\% interval & $p_\text{Holm}$ & \\
\midrule
E-ov1-en2zh & En$\rightarrow$Zh & tcpBLEU-ov, \sys{}-Simul $-$ causal relabel & 8.54 & [3.82, 14.28] & $<$0.001 & holds \\
E-ov2-en2zh & En$\rightarrow$Zh & concurrent-speaker recall, same pair & 0.168 & [0.145, 0.194] & $<$0.001 & holds \\
E-lat-en2zh & En$\rightarrow$Zh & text lag, overlapped $-$ single ($<$0.5\,s) & -0.36 & [-0.425, -0.299] & $<$0.001 & holds \\
E-pol-en2zh & En$\rightarrow$Zh & text lag, LA2 (same model) $-$ \sys{}-Simul & 0.81 & [0.553, 1.081] & $<$0.001 & holds \\
E-ov1-zh2en & Zh$\rightarrow$En & tcpBLEU-ov, \sys{}-Simul $-$ causal relabel & 8.67 & [2.63, 12.94] & 0.002 & holds \\
E-ov2-zh2en & Zh$\rightarrow$En & concurrent-speaker recall, same pair & 0.258 & [0.219, 0.302] & $<$0.001 & holds \\
E-lat-zh2en & Zh$\rightarrow$En & text lag, overlapped $-$ single ($<$0.5\,s) & -0.21 & [-0.264, -0.167] & $<$0.001 & holds \\
E-pol-zh2en & Zh$\rightarrow$En & text lag, LA2 (same model) $-$ \sys{}-Simul & 0.23 & [0.107, 0.374] & $<$0.001 & holds \\
E-ov1-ali & AliMeeting & tcpBLEU-ov, \sys{}-Simul $-$ causal relabel & 0.63 & [-2.43, 3.51] & 0.381 & tie / fails \\
E-ov2-ali & AliMeeting & concurrent-speaker recall, same pair & 0.172 & [0.131, 0.205] & $<$0.001 & holds \\
E-lat-ali & AliMeeting & text lag, overlapped $-$ single ($<$0.5\,s) & -0.23 & [-0.325, -0.119] & $<$0.001 & holds \\
E-pol-ali & AliMeeting & text lag, LA2 (same model) $-$ \sys{}-Simul & 1.89 & [1.327, 2.496] & $<$0.001 & holds \\
\bottomrule
\end{tabular}
\caption{The twelve confirmatory tests of the overlap study, Holm-corrected together. A test holds when $p_\text{Holm}<0.05$ and
the interval excludes 0 (E-lat: lies below the 0.5\,s bound). Exact sign-flip permutation over meetings; intervals from
10,000 meeting resamples, BLEU differences with a decoder-noise draw per resample.}
\label{tab:reg8}
\end{table*}

